\documentclass[10pt,a4paper]{amsart}
\usepackage{amsmath,amssymb,amsthm,graphicx}
\usepackage{verbatim}
\newenvironment{reference}{\comment}{\endcomment}
%\newenvironment{reference}{}{}
\newenvironment{slogan}{\comment}{\endcomment}
\newenvironment{history}{\comment}{\endcomment}
% For commutative diagrams we use Xy-pic
\usepackage[all]{xy}

% We use 2cell for 2-commutative diagrams.
\xyoption{2cell}
% Theorem environments.
%
\theoremstyle{plain}
\newtheorem{theorem}[subsection]{Theorem}
\newtheorem{proposition}[subsection]{Proposition}
\newtheorem{lemma}[subsection]{Lemma}

\theoremstyle{definition}
\newtheorem{definition}[subsection]{Definition}
\newtheorem{example}[subsection]{Example}
\newtheorem{exercise}[subsection]{Exercise}
\newtheorem{situation}[subsection]{Situation}

\theoremstyle{remark}
\newtheorem{remark}[subsection]{Remark}
\newtheorem{remarks}[subsection]{Remarks}

\numberwithin{equation}{subsection}

% Macros
%
\def\lim{\mathop{\mathrm{lim}}\nolimits}
\def\colim{\mathop{\mathrm{colim}}\nolimits}
\def\Spec{\mathop{\mathrm{Spec}}}
\def\Hom{\mathop{\mathrm{Hom}}\nolimits}
\def\Ext{\mathop{\mathrm{Ext}}\nolimits}
\def\SheafHom{\mathop{\mathcal{H}\!\mathit{om}}\nolimits}
\def\SheafExt{\mathop{\mathcal{E}\!\mathit{xt}}\nolimits}
\def\Sch{\mathit{Sch}}
\def\Mor{\mathop{\mathrm{Mor}}\nolimits}
\def\Ob{\mathop{\mathrm{Ob}}\nolimits}
\def\Sh{\mathop{\mathit{Sh}}\nolimits}
\def\NL{\mathop{N\!L}\nolimits}
\def\CH{\mathop{\mathrm{CH}}\nolimits}
\def\proetale{{pro\text{-}\acute{e}tale}}
\def\etale{{\acute{e}tale}}
\def\QCoh{\mathit{QCoh}}
\def\Ker{\mathop{\mathrm{Ker}}}
\def\Im{\mathop{\mathrm{Im}}}
\def\Coker{\mathop{\mathrm{Coker}}}
\def\Coim{\mathop{\mathrm{Coim}}}

% Boxtimes
%
\DeclareMathSymbol{\boxtimes}{\mathbin}{AMSa}{"02}

%
% Macros for moduli stacks/spaces
%
\def\QCohstack{\mathcal{QC}\!\mathit{oh}}
\def\Cohstack{\mathcal{C}\!\mathit{oh}}
\def\Spacesstack{\mathcal{S}\!\mathit{paces}}
\def\Quotfunctor{\mathrm{Quot}}
\def\Hilbfunctor{\mathrm{Hilb}}
\def\Curvesstack{\mathcal{C}\!\mathit{urves}}
\def\Polarizedstack{\mathcal{P}\!\mathit{olarized}}
\def\Complexesstack{\mathcal{C}\!\mathit{omplexes}}
% \Pic is the operator that assigns to X its picard group, usage \Pic(X)
% \Picardstack_{X/B} denotes the Picard stack of X over B
% \Picardfunctor_{X/B} denotes the Picard functor of X over B
\def\Pic{\mathop{\mathrm{Pic}}\nolimits}
\def\Picardstack{\mathcal{P}\!\mathit{ic}}
\def\Picardfunctor{\mathrm{Pic}}
\def\Deformationcategory{\mathcal{D}\!\mathit{ef}}
\usepackage{cleveref}
\crefname{originallemma}{original lemma}{original lemmas}
\Crefname{originallemma}{Original lemma}{Original lemmas}
\makeatletter
\newlabel{lemma-monomorphism-cartesian-square-inertia}{{lemma-monomorphism-cartesian-square-inertia}{source}{}{sourcepointer-1}{}}
\newlabel{lemma-monomorphism-cartesian-square-inertia@cref}{{[originallemma][0][]lemma-monomorphism-cartesian-square-inertia}{[1][0][]source}{}{}{}}
\makeatother
\crefname{originallemma}{original lemma}{original lemmas}
\Crefname{originallemma}{Original lemma}{Original lemmas}
\makeatletter
\newlabel{lemma-representable-decent}{{lemma-representable-decent}{source}{}{sourcepointer-2}{}}
\newlabel{lemma-representable-decent@cref}{{[originallemma][0][]lemma-representable-decent}{[1][0][]source}{}{}{}}
\makeatother
\crefname{originalproposition}{original proposition}{original propositions}
\Crefname{originalproposition}{Original proposition}{Original propositions}
\makeatletter
\newlabel{proposition-decent-sober}{{proposition-decent-sober}{source}{}{sourcepointer-3}{}}
\newlabel{proposition-decent-sober@cref}{{[originalproposition][0][]proposition-decent-sober}{[1][0][]source}{}{}{}}
\makeatother
\crefname{originalreference}{original reference}{original references}
\Crefname{originalreference}{Original reference}{Original references}
\makeatletter
\newlabel{stacks-properties-definition-reduced-induced-stack}{{stacks-properties-definition-reduced-induced-stack}{source}{}{sourcepointer-4}{}}
\newlabel{stacks-properties-definition-reduced-induced-stack@cref}{{[originalreference][0][]stacks-properties-definition-reduced-induced-stack}{[1][0][]source}{}{}{}}
\makeatother
\title{Points, irreducible closed subsets and integral closed substacks of a decent stack with quasi-compact inertia correspond}
\author{The Stacks Project Authors}
\date{stacks-project:stacks-morphisms.tex:local-snapshot-20261012; GNU FDL 1.2 or later}
\begin{document}
\maketitle
\noindent\textit{Editorial note.} All delivered mathematical text is the author's own, copied from the pinned Stacks Project chapter source sources/stacks-algebraic-stacks-morphisms/pinned-original-source.tex; no statement or proof step was rewritten, reordered or removed.\par\medskip
\noindent\textit{Editorial note.} Preamble provenance: the Stacks Project keeps one single shared preamble for every chapter; this card takes it directly from the staged sources/stacks-algebraic-stacks-morphisms/preamble.tex as three verbatim spans (lines 8-12, 14-18 and 154-217), so no derived file was written and nothing was added or redefined. The wrapper class and the base packages are supplied by the assembler.\par\medskip
\noindent\textit{Editorial note.} Bibliography provenance: neither the counted statement nor its proof contains a citation (this lane rescans the whole span before writing the card), so the card carries no bibliography and no entry was reconstructed.\par\medskip
\noindent\textit{Editorial note.} Selection provenance: this card was opened under the file-based control contract (CONTROL.md), with the mechanical pre-card gate run on its statement and proof interval before the card was written.\par\medskip
\noindent\textit{Allowed standard prerequisites (editorial).}\begin{itemize}
\item The statement that the space of a decent algebraic stack is sober. Invoked at the start of the counted proof (source lines 10722-10723). Gives the bijection between points and irreducible closed subsets.
\item The construction of the reduced induced closed substack structure on a closed subset. Invoked inside the counted proof (source lines 10725-10728) as a cross-chapter pointer. Produces the closed substack with a prescribed underlying closed subset.
\item The permanence of decency for a representable morphism. Invoked inside the counted proof (source lines 10731-10732). Makes the reduced closed substack decent.
\item The base change description of the inertia of a closed substack. Invoked near the end of the counted proof (source lines 10733-10735). Gives quasi-compactness of the inertia of the substack.
\end{itemize}\medskip
\ifdefined\hypertarget\hypertarget{sourcepointer-1}{}\fi
\noindent\textit{Source reference.} Original lemma lemma-monomorphism-cartesian-square-inertia: sources/stacks-algebraic-stacks-morphisms/pinned-original-source.tex, source lines 1480--1480. Navigation pointer only: the counted text refers to the Stacks item lemma-monomorphism-cartesian-square-inertia, which is an existing labeled result of this same chapter and is not part of this card. The pointer carries the target's own label as its number and its definition line as its span; its content is not reproduced here.\par\medskip
\ifdefined\hypertarget\hypertarget{sourcepointer-2}{}\fi
\noindent\textit{Source reference.} Original lemma lemma-representable-decent: sources/stacks-algebraic-stacks-morphisms/pinned-original-source.tex, source lines 10425--10425. Navigation pointer only: the counted text refers to the Stacks item lemma-representable-decent, which is an existing labeled result of this same chapter and is not part of this card. The pointer carries the target's own label as its number and its definition line as its span; its content is not reproduced here.\par\medskip
\ifdefined\hypertarget\hypertarget{sourcepointer-3}{}\fi
\noindent\textit{Source reference.} Original proposition proposition-decent-sober: sources/stacks-algebraic-stacks-morphisms/pinned-original-source.tex, source lines 10575--10575. Navigation pointer only: the counted text refers to the Stacks item proposition-decent-sober, which is an existing labeled result of this same chapter and is not part of this card. The pointer carries the target's own label as its number and its definition line as its span; its content is not reproduced here.\par\medskip
\ifdefined\hypertarget\hypertarget{sourcepointer-4}{}\fi
\noindent\textit{Source reference.} Original reference stacks-properties-definition-reduced-induced-stack: sources/stacks-algebraic-stacks-morphisms/pinned-original-source.tex, source lines 6968--6968. Navigation pointer only: the counted text refers to the Stacks item stacks-properties-definition-reduced-induced-stack, which belongs to another chapter of the project and is not part of this card. The pointer carries the target's own label as its number because this chapter's staged snapshot does not contain the other chapter's tag table; its content is not reproduced here.\par\medskip
\begin{lemma}
\label{lemma-decent-irreducible-closed}
Let $\mathcal{X}$ be a decent algebraic stack such that
$\mathcal{I}_\mathcal{X} \to \mathcal{X}$ is quasi-compact.
There are canonical bijections between the following sets:
\begin{enumerate}
\item the set of points of $\mathcal{X}$, i.e., $|\mathcal{X}|$,
\item the set of irreducible closed subsets of $|\mathcal{X}|$,
\item the set of integral closed substacks of $\mathcal{X}$.
\end{enumerate}
The bijection from (1) to (2) sends $x$ to $\overline{\{x\}}$.
The bijection from (3) to (2) sends $\mathcal{Z}$ to $|\mathcal{Z}|$.
\end{lemma}

\begin{proof}
Our map defines a bijection between (1) and (2) as $|\mathcal{X}|$ is
sober by  Proposition \ref{proposition-decent-sober}.
Given $T \subset |\mathcal{X}|$ closed and irreducible, there is a
unique reduced closed substack $\mathcal{Z} \subset \mathcal{X}$
such that $|\mathcal{Z}| = T$, namely, $\mathcal{Z}$
is the reduced induced subspace structure on $T$, see
Properties of Stacks, Definition
\ref{stacks-properties-definition-reduced-induced-stack}.
Then $\mathcal{Z}$ is an integral algebraic stack because it is
decent (Lemma \ref{lemma-representable-decent}), the
morphism $\mathcal{I}_\mathcal{Z} \to \mathcal{Z}$ is quasi-compact
(as the base change of $\mathcal{I}_\mathcal{X} \to \mathcal{X}$, see
Lemma \ref{lemma-monomorphism-cartesian-square-inertia}),
$\mathcal{Z}$ is reduced, and $|\mathcal{Z}|$ is irreducible.
\end{proof}

\end{document}
